\input{sol_pre.tex}
\input{q_2026-02-28_1.tex}
\begin{document}
\input{sol_head.tex}

%================= 1
\Savol{1}
$R=\underbrace{11\ldots1}_{50}$: \ $6R+5R=11R=10R+R=\underbrace{11\ldots1}_{50}0+\underbrace{11\ldots1}_{50}$.\\
Oxirgi raqam 1, birinçi raqam 1, oraliqdagi 49 ta raqam $1+1=2$: \ $1\underbrace{22\ldots2}_{49}1$.
\Javob{A}

%================= 2
\Savol{2}
$(2{,}28-4{,}35):2{,}25-(0{,}5-1{,}875)=-2{,}07:2{,}25+1{,}375=-0{,}92+1{,}375=0{,}455$.
\Javob{C) 0,455}

%================= 3
\Savol{3}
Uçraşgunça $t$ soat: velosipedçi $t-\frac14$ soat yurgan.
\[50t+20\left(t-\tfrac14\right)=275\ \Rightarrow\ 70t=280\ \Rightarrow\ t=4,\qquad 50\cdot4=200\ \text{km}.\]
\Javob{B) 200}

%================= 4
\Savol{4}
Katetlar $a$, $b$: \ $a^2+b^2=45$; \ $a\left(1+\frac43\right)+b\left(1+\frac16\right)=14\Rightarrow\frac73a+\frac76b=14\Rightarrow2a+b=12$.\\
$a^2+(12-2a)^2=45\Rightarrow5a^2-48a+99=0\Rightarrow a=3$ ($a=6{,}6$ da $b<0$), \ $b=6$. \ $S=\frac{3\cdot6}2=9$.
\Javob{A) 9}

%================= 5
\Savol{5}
$\left(\frac23\right)^{-3}=\frac{27}8$, \ $1{,}875=\frac{15}8$: \quad $\frac{27}8\cdot\frac8{15}=\frac95=1\frac45$.
\Javob{C) $1\frac45$}

%================= 6
\Savol{6}
$12345654321=111111^2$, \ $666666=6\cdot111111$: \quad $\sqrt{\left(\frac{6\cdot111111}{111111}\right)^2}=6$.
\Javob{B) 6}

%================= 7
\Savol{7}
$7-\sqrt{40}=(\sqrt5-\sqrt2)^2$ $\Rightarrow$ $\sqrt[6]{7-\sqrt{40}}=\sqrt[3]{\sqrt5-\sqrt2}$:
\[\sqrt[3]{(\sqrt5-\sqrt2)(\sqrt5+\sqrt2)}\cdot\sqrt[3]9=\sqrt[3]{27}=3.\]
\Javob{B) 3}

%================= 8
\Savol{8}
$a_1+2d=2a_1+d\Rightarrow d=a_1$; \ $a_1\cdot2a_1\cdot3a_1=6a_1^3=384\Rightarrow a_1=4$.\\
$S_{10}=\frac{2\cdot4+9\cdot4}2\cdot10=220$.
\Javob{C) 220}

%================= 9
\Savol{9}
$b_1(1+q+q^2)=19$, \ $b_1^2(1+q^2+q^4)=133$, \ $1+q^2+q^4=(1+q+q^2)(1-q+q^2)$ $\Rightarrow$ $b_1(1-q+q^2)=7$.\\
$2b_1q=12\Rightarrow b_2=6$; \ $b_1+b_3=13$, $b_1b_3=36$, kamayuvçi: $b_1=9$, $q=\frac23$.
\[S_5=\frac{9\left(1-\frac{32}{243}\right)}{\frac13}=\frac{211}9=23\frac49.\]
\Javob{B) $23\frac49$}

%================= 10
\Savol{10}
$x=4y$: \quad $\dfrac{5y+3}{10y+6}=\dfrac{5y+3}{2(5y+3)}=\dfrac12$.
\Javob{B) $\frac12$}

%================= 11
\Savol{11}
Surat: $\dfrac{y+x+1}{xy}(x+y-1)=\dfrac{(x+y)^2-1}{xy}$; \ maxraj: $\dfrac{y^2+x^2+2xy-1}{x^2y^2}=\dfrac{(x+y)^2-1}{x^2y^2}$.
\[\frac{(x+y)^2-1}{xy}\cdot\frac{x^2y^2}{(x+y)^2-1}=xy.\]
\Javob{A) $xy$}

%================= 12
\Savol{12}
$\sin\alpha=\frac45$, \ $\tg\alpha=\frac43$, \ $\ctg\alpha=\frac34$: \ $0{,}75<0{,}8<1{,}33$ $\Rightarrow$ $c<a<b$.
\Javob{C) $c<a<b$}

%================= 13
\Savol{13}
$\beta=180^\circ-2\alpha\Rightarrow\cos\beta=-\cos2\alpha$: \ $\cos\alpha-\sqrt3(2\cos^2\alpha-1)=0$.
\[2\sqrt3c^2-c-\sqrt3=0\ \Rightarrow\ c=\frac{1\pm5}{4\sqrt3}\ \Rightarrow\ \cos\alpha=\frac{\sqrt3}2\ \ (\text{ikkinçisida }\beta<0).\]
$\alpha=30^\circ$, $\beta=120^\circ$, \ $|\alpha-\beta|=90^\circ$.
\Javob{C) $90^\circ$}

%================= 14
\Savol{14}
$u=2^{x^2-x}$: \ $u^2-68u+256=0\Rightarrow u=64$ yoki $u=4$ $\Rightarrow$ $x^2-x=6$ yoki $x^2-x=2$.\\
Har birida ildizlar yiğindisi 1 (ikkalasida $D>0$, ildizlar har xil). Jami $2$.
\Javob{A) 2}

%================= 15
\Savol{15}
$\log_32\cdot\log_{32}x=\frac15\log_32\cdot\log_2x=\frac15\log_3x$, \ $\log_{\sqrt x}3=2\log_x3$. \ $t=\log_3x$:
\[\frac23-\frac1t=\frac t5-\frac4{5t}\ \Rightarrow\ 3t^2-10t+3=0\ \Rightarrow\ t=3,\ t=\frac13.\]
$x=27$, $x=\sqrt[3]3$; \ köpaytma $27\sqrt[3]3$.
\Javob{C) $27\sqrt[3]3$}

%================= 16
\Savol{16}
$\dfrac{2027!}{2025!(2026+1)}=2026$: \ $x^4-2026x^2+2025=0\Rightarrow x^2=1;\ 2025$ $\Rightarrow$ $x=\pm1,\ \pm45$ — 4 ta.
\Javob{D) 4 ta}

%================= 17
\Savol{17}
$|x|\ge45$ $\Rightarrow$ $x^4\ge45^4=2025^2$. Tenglama: $x^4-2025^2=\sqrt{|x|-45}$.\\
$|x|=45$: $0=0$ — ildiz. \ $|x|\ge46$ (butun): çap tomon $\ge46^4-45^4>\sqrt{|x|-45}$, çunki $x^4-45^4\ge(|x|-45)\cdot4\cdot45^3>\sqrt{|x|-45}$.\\
Butun yeçimlar: $x=\pm45$ — 2 ta.
\Javob{B) 2 ta}

%================= 18
\Savol{18}
Maxraj $(x-3)^2>0$ ($x\ne3$). \ Surat: $(|x|-2)(|x|-5)<0\Rightarrow2<|x|<5$.\\
Butunlar: $\pm3,\pm4$, lekin $x\ne3$: \ $-4,-3,4$ — 3 ta.
\Javob{C) 3 ta}

%================= 19
\Savol{19}
$x\le2$, \ $\sqrt{(x-2)^2}=2-x$. \ $t=\sqrt{2-x}\ge0$: $t^2-5t-6<0\Rightarrow t<6\Rightarrow2-x<36\Rightarrow x>-34$.\\
$-33\le x\le2$ — 36 ta.
\Javob{B) 36}

%================= 20
\Savol{20}
$f(x)=a(x-3)^2+5$, \ $f(0)=9a+5=\frac{16}5\Rightarrow a=-\frac15$.
\[f(18)=-\frac15\cdot225+5=-40.\]
\Javob{C) $-40$}

%================= 21
\Savol{21}
$\dfrac{x^2}{x^2+1}=1-\dfrac1{x^2+1}$: \quad $\displaystyle\int_{-1}^1=2-\arctg x\Big|_{-1}^1=2-\frac\pi2$.
\Javob{A) $2-\frac\pi2$}

%================= 22
\Savol{22}
$P(x)=\sum C_kx^{20-k}$, \ $P''(x)=\sum(20-k)(19-k)C_kx^{18-k}$. \ $x=-1$ da $(-1)^{18-k}=(-1)^k$ — aynan berilgan yiğindi ($C_{19},C_{20}$ hadlari nolga aylanadi).
\[P''(x)=3\cdot20\cdot19\cdot9(3x+4)^{18},\qquad P''(-1)=3\cdot380\cdot9\cdot1=10260.\]
\Javob{B) 10260}

%================= 23
\Savol{23}
Markaziy burçak $90^\circ$: vatar $=R\sqrt2=8\Rightarrow R=4\sqrt2$.
\Javob{C) $4\sqrt2$}

%================= 24
\Savol{24}
Funksiyada har bir $x$ ga bitta $y$ mos keladi (vertikal töğri çiziq grafikni köpi bilan bir marta kesadi).\\
1) aylana, 3) öng yarim aylana, 4) yotiq parabola — ba'zi $x$ larda ikki nuqta. Faqat 2) — funksiya.
\Javob{A) 1}

%================= 25
\Savol{25}
Bissektrisa xossasi: $\frac{AB}{BC}=\frac{\sqrt5}{2\sqrt5}=\frac12$: \ $AB=k$, $BC=2k$, \ $AC=3\sqrt5$.
\[k^2+4k^2=45\Rightarrow k=3,\qquad S=\frac12\cdot3\cdot6=9.\]
\Javob{B) 9}

%================= 26
\Savol{26}
$AC=12$. $BE$ bissektrisa: $\frac{BC}{AB}=\frac{EC}{AE}=2\Rightarrow BC=12$. \ $AD$ bissektrisa: $\frac{BD}{DC}=\frac{AB}{AC}=\frac12\Rightarrow BD=4$, $DC=8$.\\
$K$ — ikki bissektrisa kesişmasi, ya'ni içki aylana markazi.
\fig{\begin{tikzpicture}[scale=0.42]
\coordinate (A) at (0,0); \coordinate (C) at (12,0); \coordinate (B) at (1.5,5.809); \coordinate (E) at (4,0); \coordinate (D) at (5,3.873); \coordinate (K) at (3,2.324);
\fill[blue!15] (K)--(D)--(C)--(E)--cycle;
\draw (A)--(B)--(C)--cycle; \draw (A)--(D); \draw (B)--(E);
\draw[axs] (K)--(3,0); \draw[ax] (C)--(K);
\node[below left] at (A){$A$}; \node[below right] at (C){$C$}; \node[above] at (B){$B$}; \node[below] at (E){$E$}; \node[above right] at (D){$D$}; \node[left] at (K){$K$};
\node[aux,font=\scriptsize] at (3.4,1){$r$};
\end{tikzpicture}}
$\triangle ABC$ ($AC=BC=12$, $AB=6$): balandlik $\sqrt{144-9}=3\sqrt{15}$, \ $S=9\sqrt{15}$, \ $p=15$, \ $r=\frac{3\sqrt{15}}5$.
\[S_{ABD}=\tfrac4{12}S=3\sqrt{15},\quad S_{ABE}=\tfrac4{12}S=3\sqrt{15},\quad S_{ABK}=\tfrac12\cdot6\cdot r=\tfrac{9\sqrt{15}}5.\]
\[S_{KDCE}=S-S_{ABD}-S_{ABE}+S_{ABK}=9\sqrt{15}-6\sqrt{15}+\frac{9\sqrt{15}}5=\frac{24\sqrt{15}}5.\]
\Javob{C) $\frac{24\sqrt{15}}5$}

%================= 27
\Savol{27}
$C$ dan $CE\parallel AB$: $\triangle ECD$ tomonlari $12,\ 16,\ 27-7=20$ — $12^2+16^2=20^2$, töğri burçakli.
\fig{\begin{tikzpicture}[scale=0.2]
\coordinate (A) at (0,0); \coordinate (D) at (27,0); \coordinate (B) at (7.2,9.6); \coordinate (C) at (14.2,9.6); \coordinate (E) at (7,0);
\draw (A)--(B)--(C)--(D)--cycle; \draw[axs] (C)--(E); \draw[ax] (C)--(14.2,0);
\node[below left] at (A){$A$}; \node[below right] at (D){$D$}; \node[above] at (B){$B$}; \node[above] at (C){$C$}; \node[below] at (E){$E$};
\end{tikzpicture}}
$h=\dfrac{12\cdot16}{20}=9{,}6$, \quad $S=\dfrac{7+27}2\cdot9{,}6=163{,}2=163\frac15$.
\Javob{A) $163\frac15$}

%================= 28
\Savol{28}
Bir uçdan çiqqan ikki eng kiçik diagonal (qöşni-qöşni uçlarga) orasidagi içki çizilgan burçak qolgan $n-4$ ta tomon yoyiga tiraladi:
\[\frac{(n-4)\cdot360^\circ}{2n}=120^\circ\ \Rightarrow\ 180n-720=120n\ \Rightarrow\ n=12.\]
\Javob{B) 12}

%================= 29
\Savol{29}
$AB=\sqrt{20}=2\sqrt5$, \ $AC=\sqrt{10}$, \ $BC=\sqrt{10}$. \ $\frac{BL}{LC}=\frac{AB}{AC}=\sqrt2$ $\Rightarrow$ $BL\cdot LC=\dfrac{\sqrt2\cdot10}{(1+\sqrt2)^2}=30\sqrt2-40$.
\[AL^2=AB\cdot AC-BL\cdot LC=10\sqrt2-(30\sqrt2-40)=40-20\sqrt2.\]
\Javob{B) $\sqrt{40-20\sqrt2}$}

%================= 30
\Savol{30}
$S_{toʻla}=3S_{asos}\Rightarrow S_{yon}=2S_{asos}$. \ $S_{asos}=S_{yon}\cos\varphi\Rightarrow\cos\varphi=\frac12$, \ $\varphi=60^\circ$, \ $\tg\varphi=\sqrt3$.
\Javob{C) $\sqrt3$}

%================= 31
\Savol{31}
$A\cap B$ — $\text{EKUB}(60;48)=12$ ning böluvçilari: $1,2,3,4,6,12$ — 6 ta. \ Qism töplamlar: $2^6=64$.
\Javob{C) 64}

%================= 32
\Savol{32}
$\overline{abc}$, $a\ge1$, $a+b+c=7$: \ $a'=a-1$, $a'+b+c=6$ $\Rightarrow$ $C_8^2=28$ ta son.\\
Juft: $c=0$: $a+b=7$ — 7 ta; $c=2$: 5 ta; $c=4$: 3 ta; $c=6$: 1 ta; jami 16. \ $P=\frac{16}{28}=\frac47$.
\Javob{D) $\frac47$}

%================= 33-35
\newpage
\Savol{33}
Konus $R$, $H$; silindr $r$, $h$. Öq kesimida öxşaşlik: $\dfrac hH=\dfrac{R-r}R$.
\fig{\begin{tikzpicture}[scale=0.6]
\draw (-3,0)--(3,0)--(0,4.5)--cycle; \draw[okgreen,thick] (-2,0) rectangle (2,1.5);
\draw[ax] (0,0)--(0,4.5); \node[below] at (1,0){$r$}; \node[below] at (2.6,0){}; \node[right] at (2,0.75){$h$};
\node[left,font=\scriptsize] at (0,3){$H$}; \node[below] at (-1.5,0){$R$};
\end{tikzpicture}}
\[V_{sil}=\pi r^2H\left(1-\frac rR\right),\qquad V'=\pi H\left(2r-\frac{3r^2}R\right)=0\ \Rightarrow\ r=\frac23R.\]
\textbf{33.} $\frac rR=\frac23$. \quad \textbf{34.} $\frac hH=1-\frac23=\frac13$. \quad \textbf{35.} $\dfrac{\pi\cdot\frac49R^2\cdot\frac H3}{\frac13\pi R^2H}=\frac49$.
\Javob{33 — B; \quad 34 — A; \quad 35 — E}

%================= 36
\newpage
\Savol{36}
Birinçisini 3 ga, ikkinçisini 5 ga köpaytirib ayiramiz (öng tomonlar 30):
\[28x^2+36xy+11y^2=0\iff(2x+y)(14x+11y)=0.\]
1) $y=-2x$: \ $10x^2-36x^2+32x^2=6x^2=6\Rightarrow x=\pm1$: \ $(1;-2)$, $(-1;2)$.

2) $y=-\frac{14}{11}x$: \ $10x^2-\frac{252}{11}x^2+\frac{1568}{121}x^2=\frac{6x^2}{121}=6\Rightarrow x=\pm11$: \ $(11;-14)$, $(-11;14)$.

a) 4 ta yeçim. \qquad b) $x+y$: $-1,\ 1,\ -3,\ 3$ $\Rightarrow$ eng kattasi 3.
\Javob{a) 4; \quad b) 3}

%================= 37
\newpage
\Savol{37}
$a=\sin x$, $b=\sin2x$, $c=\sin3x$: \ $(a+b+c)^3-a^3-b^3-c^3=3(a+b)(b+c)(c+a)=0$.
\[a+b=2\sin\tfrac{3x}2\cos\tfrac x2,\quad b+c=2\sin\tfrac{5x}2\cos\tfrac x2,\quad c+a=2\sin2x\cos x.\]
Ildizlar: $x=\frac{2\pi k}3$, \ $x=\pi+2\pi k$, \ $x=\frac{2\pi k}5$, \ $x=\frac{\pi k}2$.

a) Eng kiçik musbat: $\frac{2\pi}5$ \ ($<\frac\pi2<\frac{2\pi}3$).

b) $[-\pi;\pi]$ da: $0,\ \pm\frac{2\pi}5,\ \pm\frac\pi2,\ \pm\frac{2\pi}3,\ \pm\frac{4\pi}5,\ \pm\pi$ — 11 ta.
\Javob{a) $\frac{2\pi}5$; \quad b) 11}

%================= 38
\newpage
\Savol{38}
$f(3)=0\Rightarrow3a+b=0$. \ $E(y)$ dan çiqarilgan son — gorizontal asimptota: $\frac ac=2$.\\
$f=f^{-1}$: teskari funksiya $f^{-1}(x)=\dfrac{-dx+b}{cx-a}$, tenglik $\Rightarrow$ $d=-a$.\\
$c=1$ deb: $a=2$, $b=-6$, $d=-2$: \ $f(x)=\dfrac{2x-6}{x-2}$.

a) $x-2\ne0$: aniqlaniş sohasiga $x=2$ tegişli emas.

b) $f^{-1}(6)=f(6)=\dfrac{12-6}{6-2}=\dfrac32$.
\Javob{a) 2; \quad b) $\frac32$}

%================= 39
\newpage
\Savol{39}
$f(x)=\frac18(x-6)^2-\frac92$ — uçi $V(6;-4{,}5)$. \ Aylana: $(x-9)^2+(y-6)^2=20$ — markaz $M(9;6)$, $r=2\sqrt5$.

a) $d_1=MV=\sqrt{3^2+10{,}5^2}=\sqrt{\frac{477}4}=\frac{3\sqrt{53}}2$.

b) Eng qisqa masofa parabola normalining $M$ dan ötgan nuqtasida. $P(x_0;f(x_0))$ da urinma burçak koeffitsiyenti $f'(x_0)=\frac{x_0}4-\frac32$; $\overrightarrow{PM}\perp$ urinma:
\[(9-x_0)+(6-f(x_0))f'(x_0)=0.\]
$x_0=14$: $f=3{,}5$, $f'=2$: \ $-5+2{,}5\cdot2=0$ — mos. \ $PM=\sqrt{25+6{,}25}=\frac{5\sqrt5}2$.
\[d_2=PM-r=\frac{5\sqrt5}2-2\sqrt5=\frac{\sqrt5}2.\]
\Javob{a) $\frac{3\sqrt{53}}2$; \quad b) $\frac{\sqrt5}2$}

%================= 40
\newpage
\Savol{40}
$f-g=-3x^2+(4-p)x-q$ — parabola; soha uning ildizlari $x_1,x_2$ orasida va parabola öqiga nisbatan simmetrik. Teng ikkiga böluvçi vertikal — simmetriya öqi: $\frac{x_1+x_2}2=\frac43$.

a) $x_1+x_2=\frac83$.

b) $S=\displaystyle\int_{x_1}^{x_2}3(x-x_1)(x_2-x)dx=\frac{3(x_2-x_1)^3}6=\frac{32}{27}\Rightarrow x_2-x_1=\frac43$ $\Rightarrow$ $x_1=\frac23$, $x_2=2$.\\
$x_1+x_2=\frac{4-p}3=\frac83\Rightarrow p=-4$; \ $x_1x_2=\frac q3=\frac43\Rightarrow q=4$. \ $pq=-16$.
\Javob{a) $\frac83$; \quad b) $-16$}

%================= 41
\newpage
\Savol{41}
$O$ — bissektrisalar kesişmasi. $EF\parallel BC$ $\Rightarrow$ $\angle EOB=\angle OBC=\angle OBE$ $\Rightarrow$ $EO=EB=8$; xuddi şunday $FO=FC=12$, \ $EF=20$.
\fig{\begin{tikzpicture}[scale=0.2]
\coordinate (B) at (0,0); \coordinate (C) at (30,0); \coordinate (A) at (3,23.81); \coordinate (O) at (9,7.937);
\coordinate (E) at (1,7.937); \coordinate (F) at (21,7.937); \coordinate (D) at (12,0);
\draw (A)--(B)--(C)--cycle; \draw (O) circle (7.937); \draw (E)--(F); \draw (A)--(D);
\draw[axs] (B)--(O); \draw[axs] (C)--(O);
\node[left] at (A){$A$}; \node[below left] at (B){$B$}; \node[below right] at (C){$C$}; \node[left] at (E){$E$}; \node[right] at (F){$F$}; \node[below] at (D){$D$}; \node[above right] at (O){$O$};
\end{tikzpicture}}
$\triangle AEF\sim\triangle ABC$, \ $k=\frac{EF}{BC}=\frac23$ $\Rightarrow$ $EB=\frac13AB=8$, $FC=\frac13AC=12$ $\Rightarrow$ $AB=24$, $AC=36$.

a) $\dfrac{BD}{DC}=\dfrac{AB}{AC}=\dfrac23$, \ $DC=\dfrac35\cdot30=18$.

b) $|AB-AC|=12$.
\Javob{a) 18; \quad b) 12}

%================= 42
\newpage
\Savol{42}
a) Kvadrat içidagi nuqta uçun $PA^2+PC^2=PB^2+PD^2$ ($P$ dan tomonlarga perpendikulyarlar va Pifagor teoremasi): \ $1+49=25+PD^2\Rightarrow PD=5$.

b) $\triangle BPA$ ni $B$ atrofida $90^\circ$ ga buramiz: $A\to C$, $P\to P'$. \ $BP'=BP=5$, $\angle PBP'=90^\circ$ $\Rightarrow$ $PP'=5\sqrt2$; \ $CP'=AP=1$.
\fig{\begin{tikzpicture}[scale=0.62]
\coordinate (A) at (0,0); \coordinate (B) at (0,5.657); \coordinate (C) at (5.657,5.657); \coordinate (D) at (5.657,0); \coordinate (P) at (0.707,0.707); \coordinate (Q) at (4.95,6.364);
\draw (A)--(B)--(C)--(D)--cycle; \draw (P)--(A); \draw (P)--(B); \draw (P)--(C); \draw (P)--(D);
\draw[axs] (B)--(Q)--(C); \draw[ax] (P)--(Q);
\node[below left] at (A){$A$}; \node[left] at (B){$B$}; \node[right] at (C){$C$}; \node[below right] at (D){$D$}; \node[below right] at (P){$P$}; \node[above] at (Q){$P'$};
\end{tikzpicture}}
$\triangle PCP'$: $PC^2+CP'^2=49+1=50=PP'^2$ $\Rightarrow$ $\angle PCP'=90^\circ$, \ $\tg\angle PP'C=7$, $\cos=\frac1{\sqrt{50}}$, $\sin=\frac7{\sqrt{50}}$.\\
$\angle BPA=\angle BP'C=45^\circ+\angle PP'C$: \ $\cos\angle BPA=\frac{1-7}{\sqrt{50}\sqrt2}=-\frac35$.
\[AB^2=1+25-2\cdot1\cdot5\cdot\left(-\frac35\right)=32.\]
\Javob{a) 5; \quad b) 32}

%================= 43
\newpage
\Savol{43}
$DF\parallel AC$, $KF\parallel AB$. \ $\triangle GEF\sim\triangle GCK$ ($EF\parallel KC$): $\frac{EF}{KC}=\frac{EG}{GC}=\frac12$.\\
$DE=2t$, $EF=3t$ $\Rightarrow$ $KC=6t$, \ $AK=DF=5t$, \ $AC=11t$.

a) $\triangle BDE\sim\triangle BAC$: $\frac{BE}{BC}=\frac{DE}{AC}=\frac2{11}$ $\Rightarrow$ $EC=\frac9{11}BC$, \ $EG=\frac13EC=\frac3{11}BC$.

b) $\frac{AD}{AB}=\frac9{11}$, $\frac{AK}{AC}=\frac5{11}$. \ $S_{ADFK}=AD\cdot AK\sin A=2\cdot\frac9{11}\cdot\frac5{11}\cdot S_{ABC}=\frac{90}{121}\cdot242=180$.
\Javob{a) $\frac3{11}$; \quad b) 180}

%================= 44
\newpage
\Savol{44}
$O$ — markaz ($AB$ örtasi), $\rho$ — $O$ dan $CD$ çiziqqaça masofa, $d$ — $AB$ va $CD$ orasidagi masofa ($d\le\rho$), $CD=2\sqrt{9-\rho^2}$.
\[V=\frac16AB\cdot CD\cdot d\cdot\sin\varphi\le CD\cdot\rho=2\rho\sqrt{9-\rho^2}\le\rho^2+(9-\rho^2)=9.\]
Tenglik: $CD\perp AB$, $CD$ — $AB$ ga perpendikulyar va $O$ dan ötgan tekislikda, $OE=\rho=\frac3{\sqrt2}$ ($E$ — $CD$ örtasi).
\fig{\begin{tikzpicture}[scale=0.65]
\draw (0,0) circle (3);
\draw[dashed,domain=0:180,samples=40] plot ({3*cos(\x)+1.05*sin(\x)},{0.75*sin(\x)});
\draw[domain=180:360,samples=40] plot ({3*cos(\x)+1.05*sin(\x)},{0.75*sin(\x)});
\coordinate (A) at (0,3); \coordinate (B) at (0,-3); \coordinate (O) at (0,0); \coordinate (E) at (2.121,0); \coordinate (C) at (2.863,0.530); \coordinate (D) at (1.379,-0.530);
\draw (A)--(C)--(B)--(D)--cycle; \draw (C)--(D); \draw[dashed] (A)--(B);
\draw[axs] (A)--(E); \draw[ax] (O)--(E); \draw[ax] (O)--(C);
\node[above] at (A){$A$}; \node[below] at (B){$B$}; \node[left] at (O){$O$}; \node[right] at (C){$C$}; \node[below] at (D){$D$}; \node[below right,font=\scriptsize] at (E){$E$};
\end{tikzpicture}}
a) $V_{max}=9$.

b) Asos — $ABC$. $ABC$ tekislik $AB$ ni öz içiga oladi, şuning uçun $COD$ tekislikka perpendikulyar va uni $OC$ böyiça kesadi. $\triangle OEC$: $\angle E=90^\circ$, $OE=EC=\frac3{\sqrt2}$ $\Rightarrow$ $E$ dan $OC$ gaça (ya'ni $ABC$ tekislikkaça) $\frac3{\sqrt2}\sin45^\circ=\frac32$.
\[AE=\sqrt{OA^2+OE^2}=\sqrt{9+\frac92}=\frac{3\sqrt6}2,\qquad\sin\varphi=\frac{3/2}{3\sqrt6/2}=\frac1{\sqrt6}=\frac{\sqrt6}6.\]
\Javob{a) 9; \quad b) $\frac{\sqrt6}6$}

%================= 45
\newpage
\Savol{45}
$\pi r^2h=162\pi$, \ $2\pi rh=72\pi$ $\Rightarrow$ $rh=36$, $r=\frac{162}{36}=4{,}5$, $h=8$.\\
Silindrdan kesiladigan eng katta şar radiusi $\min\left(r;\frac h2\right)=\min(4{,}5;4)=4$.

a) $S=4\pi\cdot16=64\pi\approx192$ cm$^2$.

b) $V_{chiq}=162\pi-\frac43\pi\cdot64=\frac{486-256}3\pi=\frac{230\pi}3\approx230$ cm$^3$.
\Javob{a) 192 cm$^2$; \quad b) 230 cm$^3$}

\ShowJavoblar
\end{document}
